\documentclass[12pt,a4paper]{article}

% Ελληνικά
\usepackage[utf8]{inputenc}
\usepackage[greek,english]{babel}
\usepackage{alphabeta}

% Πακέτα
\usepackage{geometry}
\usepackage{graphicx}
\usepackage{booktabs}
\usepackage{longtable}
\usepackage{array}
\usepackage{multirow}
\usepackage{float}
\usepackage{caption}
\usepackage{subcaption}
\usepackage{xcolor}
\usepackage{listings}
\usepackage{amsmath}
\usepackage{amssymb}
\usepackage{hyperref}
\usepackage{fancyhdr}
\usepackage{titlesec}
\usepackage{enumitem}
\usepackage{tcolorbox}

% Ρυθμίσεις σελίδας
\geometry{
    left=2.5cm,
    right=2.5cm,
    top=2.5cm,
    bottom=2.5cm
}

% Ρυθμίσεις κώδικα Python
\definecolor{codegreen}{rgb}{0,0.6,0}
\definecolor{codegray}{rgb}{0.5,0.5,0.5}
\definecolor{codepurple}{rgb}{0.58,0,0.82}
\definecolor{backcolour}{rgb}{0.95,0.95,0.92}

\lstdefinestyle{pythonstyle}{
    backgroundcolor=\color{backcolour},
    commentstyle=\color{codegreen},
    keywordstyle=\color{blue},
    numberstyle=\tiny\color{codegray},
    stringstyle=\color{codepurple},
    basicstyle=\ttfamily\footnotesize,
    breakatwhitespace=false,
    breaklines=true,
    captionpos=b,
    keepspaces=true,
    numbers=left,
    numbersep=5pt,
    showspaces=false,
    showstringspaces=false,
    showtabs=false,
    tabsize=2,
    language=Python,
    frame=single,
    rulecolor=\color{black}
}

\lstset{style=pythonstyle}

% Ρυθμίσεις hyperref
\hypersetup{
    colorlinks=true,
    linkcolor=blue,
    filecolor=magenta,
    urlcolor=cyan,
    citecolor=green
}

% Ρυθμίσεις header/footer
\pagestyle{fancy}
\fancyhf{}
\rhead{ΑΕΜ: 6931}
\lhead{Ανάλυση Δεδομένων 2025-26}
\cfoot{\thepage}

% Custom box για σημαντικές πληροφορίες
\newtcolorbox{infobox}[1][]{
    colback=blue!5!white,
    colframe=blue!75!black,
    fonttitle=\bfseries,
    title=#1
}

\newtcolorbox{warningbox}[1][]{
    colback=orange!5!white,
    colframe=orange!75!black,
    fonttitle=\bfseries,
    title=#1
}

\newtcolorbox{successbox}[1][]{
    colback=green!5!white,
    colframe=green!75!black,
    fonttitle=\bfseries,
    title=#1
}

% Τίτλος
\title{
    \vspace{-1cm}
    \textbf{Ανάλυση Δεδομένων 2025-26}\\
    \large 1η Εργασία: Καθαρισμός και Ανάλυση Δεδομένων Ποιότητας Κρασιού\\
    \vspace{0.5cm}
    \normalsize Ερώτημα 1: Δειγματοληψία και Καθαρισμός Δεδομένων
}

\author{
    \textbf{ΑΕΜ:} 6931
}

\date{Μάρτιος 2026}

\begin{document}

\maketitle

\tableofcontents
\newpage

% ==============================================================================
\section{Εισαγωγή}
% ==============================================================================

\subsection{Σκοπός της Εργασίας}

Η παρούσα εργασία αφορά την ανάλυση ενός συνόλου δεδομένων (dataset) που περιέχει χαρακτηριστικά ποιότητας διαφόρων τύπων κρασιού. Ο στόχος του πρώτου ερωτήματος είναι:

\begin{enumerate}
    \item Η επιλογή ενός \textbf{τυχαίου δείγματος} που αποτελεί το 80\% των αρχικών δεδομένων
    \item Ο \textbf{έλεγχος ποιότητας} των δεδομένων για τον εντοπισμό προβλημάτων
    \item Ο \textbf{καθαρισμός} των δεδομένων με κατάλληλες τεχνικές
\end{enumerate}

\subsection{Τι είναι ο Καθαρισμός Δεδομένων;}

\begin{infobox}[Ορισμός: Καθαρισμός Δεδομένων (Data Cleaning)]
Ο καθαρισμός δεδομένων είναι η διαδικασία εντοπισμού και διόρθωσης (ή αφαίρεσης) λανθασμένων, ελλιπών, διπλότυπων ή ασυνεπών εγγραφών από ένα σύνολο δεδομένων. Είναι ένα κρίσιμο βήμα πριν από οποιαδήποτε ανάλυση, καθώς η ποιότητα των αποτελεσμάτων εξαρτάται άμεσα από την ποιότητα των δεδομένων.
\end{infobox}

\subsection{Γιατί είναι Σημαντικός ο Καθαρισμός Δεδομένων;}

Τα πραγματικά δεδομένα σπάνια είναι τέλεια. Συνήθη προβλήματα περιλαμβάνουν:

\begin{itemize}
    \item \textbf{Τυπογραφικά λάθη:} Λανθασμένη πληκτρολόγηση (π.χ. ``whte'' αντί ``white'')
    \item \textbf{Ελλείπουσες τιμές (Missing Values):} Κενά πεδία σε ορισμένες εγγραφές
    \item \textbf{Μη έγκυρες τιμές:} Τιμές εκτός λογικών ορίων (π.χ. pH = 999)
    \item \textbf{Διπλοεγγραφές:} Η ίδια εγγραφή εμφανίζεται πολλές φορές
\end{itemize}

Αν δεν αντιμετωπιστούν αυτά τα προβλήματα, τα αποτελέσματα της ανάλυσης θα είναι αναξιόπιστα ή παραπλανητικά.

% ==============================================================================
\section{Περιγραφή του Συνόλου Δεδομένων}
% ==============================================================================

\subsection{Πηγή και Περιεχόμενο}

Το σύνολο δεδομένων περιέχει μετρήσεις χημικών χαρακτηριστικών κρασιών (κόκκινων και λευκών) μαζί με βαθμολογίες ποιότητας. Κάθε γραμμή αντιπροσωπεύει ένα δείγμα κρασιού.

\subsection{Μεταβλητές του Dataset}

\begin{table}[H]
\centering
\caption{Περιγραφή των μεταβλητών του dataset}
\label{tab:variables}
\begin{tabular}{@{}clp{8cm}@{}}
\toprule
\textbf{Α/Α} & \textbf{Μεταβλητή} & \textbf{Περιγραφή} \\
\midrule
1 & Fixed acidity & Σταθερή οξύτητα - Μη πτητικά οξέα (g/L) \\
2 & Volatile acidity & Πτητική οξύτητα - Οξικό οξύ (g/L) \\
3 & Citric acid & Κιτρικό οξύ (g/L) \\
4 & Residual sugar & Υπολειμματικά σάκχαρα μετά τη ζύμωση (g/L) \\
5 & Chlorides & Χλωριούχα άλατα - Δείκτης αλατότητας (g/L) \\
6 & Free sulfur dioxide & Ελεύθερο SO\textsubscript{2} - Προστασία από οξείδωση (mg/L) \\
7 & Total sulfur dioxide & Συνολικό SO\textsubscript{2} (ελεύθερο + δεσμευμένο) (mg/L) \\
8 & Density & Πυκνότητα (g/cm\textsuperscript{3}) \\
9 & pH & Μέτρο της οξύτητας (κλίμακα 0-14) \\
10 & Sulphates & Θειικά άλατα (g/L) \\
11 & Alcohol & Περιεκτικότητα σε αιθανόλη (\% vol) \\
12 & Quality & Βαθμολογία ποιότητας (κλίμακα 0-10) \\
13 & Wine\_type & Τύπος κρασιού: red (κόκκινο) ή white (λευκό) \\
\bottomrule
\end{tabular}
\end{table}

\begin{infobox}[Τύποι Μεταβλητών]
\begin{itemize}
    \item \textbf{Συνεχείς (Continuous):} Μεταβλητές 1-11 - μπορούν να πάρουν οποιαδήποτε αριθμητική τιμή σε ένα εύρος
    \item \textbf{Διακριτές/Κατηγορικές:} Quality (διακριτή κλίμακα 0-10), Wine\_type (κατηγορική με τιμές ``red'' ή ``white'')
\end{itemize}
\end{infobox}

\subsection{Αρχικές Διαστάσεις}

Το αρχικό dataset περιείχε:
\begin{itemize}
    \item \textbf{8.385 γραμμές} (εγγραφές/δείγματα κρασιού)
    \item \textbf{15 στήλες} (13 μεταβλητές + 2 κενές στήλες)
\end{itemize}

% ==============================================================================
\section{Δειγματοληψία (Sampling)}
% ==============================================================================

\subsection{Τι είναι η Δειγματοληψία;}

\begin{infobox}[Ορισμός: Δειγματοληψία]
Η δειγματοληψία είναι η διαδικασία επιλογής ενός υποσυνόλου (δείγματος) από ένα μεγαλύτερο σύνολο δεδομένων (πληθυσμό). Χρησιμοποιείται για να μειώσει τον όγκο των δεδομένων προς ανάλυση ή για να δημιουργήσει διαφορετικά σύνολα για εκπαίδευση και αξιολόγηση μοντέλων.
\end{infobox}

\subsection{Γιατί τυχαία δειγματοληψία;}

Η \textbf{τυχαία δειγματοληψία} εξασφαλίζει ότι:
\begin{enumerate}
    \item Κάθε εγγραφή έχει \textbf{ίση πιθανότητα} να επιλεγεί
    \item Το δείγμα είναι \textbf{αντιπροσωπευτικό} του συνόλου
    \item Αποφεύγεται η \textbf{μεροληψία} (bias) στην επιλογή
\end{enumerate}

\subsection{Ο Ρόλος του Seed (Σπόρος)}

\begin{warningbox}[Σημαντικό: Random Seed]
Οι υπολογιστές δημιουργούν ``ψευδοτυχαίους'' αριθμούς χρησιμοποιώντας αλγορίθμους. Ο \textbf{seed} (σπόρος) είναι μια αρχική τιμή που καθορίζει την ακολουθία των τυχαίων αριθμών.

\textbf{Γιατί είναι σημαντικό:} Χρησιμοποιώντας τον ίδιο seed, παίρνουμε πάντα τα ίδια ``τυχαία'' αποτελέσματα. Αυτό επιτρέπει την \textbf{αναπαραγωγιμότητα} (reproducibility) της ανάλυσης.

Στην εργασία, χρησιμοποιούμε ως seed τον ΑΕΜ (6931), ώστε κάθε φοιτητής να εργάζεται με διαφορετικό δείγμα.
\end{warningbox}

\subsection{Υλοποίηση στην Python}

\begin{lstlisting}[caption={Κώδικας δειγματοληψίας}]
import pandas as pd

# Φόρτωση δεδομένων
df_original = pd.read_csv('Data Analysis_2026 1st Case_Data.csv')

# Ορισμός παραμέτρων
AEM = 6931              # Ο σπόρος (seed)
SAMPLE_FRACTION = 0.80  # Ποσοστό δείγματος (80%)

# Δειγματοληψία
df = df_original.sample(frac=SAMPLE_FRACTION, random_state=AEM)

# Επαναφορά index
df = df.reset_index(drop=True)
\end{lstlisting}

\textbf{Επεξήγηση του κώδικα:}
\begin{itemize}
    \item \texttt{sample()}: Η συνάρτηση της pandas για τυχαία επιλογή γραμμών
    \item \texttt{frac=0.80}: Επιλέγουμε το 80\% των γραμμών
    \item \texttt{random\_state=AEM}: Ορίζουμε τον seed για αναπαραγωγιμότητα
    \item \texttt{reset\_index()}: Επαναριθμούμε τις γραμμές από 0
\end{itemize}

\subsection{Αποτέλεσμα Δειγματοληψίας}

\begin{table}[H]
\centering
\caption{Αποτελέσματα δειγματοληψίας}
\begin{tabular}{@{}lcc@{}}
\toprule
\textbf{Μέγεθος} & \textbf{Γραμμές} & \textbf{Ποσοστό} \\
\midrule
Αρχικό dataset & 8.385 & 100\% \\
Δείγμα (ΑΕΜ: 6931) & 6.708 & 80\% \\
\bottomrule
\end{tabular}
\end{table}

% ==============================================================================
\section{Εντοπισμός Προβλημάτων στα Δεδομένα}
% ==============================================================================

Μετά τη δειγματοληψία, πραγματοποιήθηκε συστηματικός έλεγχος για τον εντοπισμό προβλημάτων. Τα προβλήματα οργανώθηκαν σε τέσσερις κατηγορίες.

\subsection{Κατηγορία 1: Τυπογραφικά Λάθη στη Στήλη wine\_type}

\subsubsection{Τι ελέγξαμε}

Η στήλη \texttt{wine\_type} πρέπει να περιέχει μόνο δύο τιμές: ``red'' ή ``white''. Ελέγξαμε τις μοναδικές τιμές που εμφανίζονται:

\begin{lstlisting}[caption={Έλεγχος μοναδικών τιμών}]
# Εμφάνιση όλων των μοναδικών τιμών
print(df['wine_type'].unique())
\end{lstlisting}

\subsubsection{Τι βρήκαμε}

\begin{table}[H]
\centering
\caption{Τυπογραφικά λάθη στη στήλη wine\_type}
\label{tab:typos}
\begin{tabular}{@{}lcc@{}}
\toprule
\textbf{Λανθασμένη Τιμή} & \textbf{Πλήθος Εμφανίσεων} & \textbf{Πιθανή Σωστή Τιμή} \\
\midrule
Redd  & 75 & red \\
whit & 90 & white \\
r3d & 64 & red \\
WHITTE & 77 & white \\
whate & 39 & white \\
999 & 59 & Μη έγκυρη \\
9999 & 45 & Μη έγκυρη \\
-999 & 64 & Μη έγκυρη \\
? & 47 & Μη έγκυρη \\
\midrule
\textbf{Σύνολο} & \textbf{560} & \\
\bottomrule
\end{tabular}
\end{table}

\begin{infobox}[Ανάλυση των λαθών]
\begin{itemize}
    \item \textbf{``Redd '', ``r3d''}: Παραλλαγές του ``red'' με τυπογραφικά λάθη ή κενά
    \item \textbf{``whit'', ``WHITTE'', ``whate''}: Παραλλαγές του ``white''
    \item \textbf{``999'', ``9999'', ``-999'', ``?''}: Κωδικοί που πιθανώς χρησιμοποιήθηκαν για να δηλώσουν ελλείπουσες ή άγνωστες τιμές
\end{itemize}
\end{infobox}

\subsection{Κατηγορία 2: Μη Αριθμητικές Τιμές σε Αριθμητικές Στήλες}

\subsubsection{Τι ελέγξαμε}

Οι 12 μεταβλητές (εκτός του wine\_type) πρέπει να περιέχουν αριθμητικές τιμές. Ελέγξαμε αν υπάρχουν τιμές που δεν μπορούν να μετατραπούν σε αριθμούς.

\begin{lstlisting}[caption={Μετατροπή σε αριθμητικές τιμές}]
# Μετατροπή με errors='coerce' μετατρέπει τα μη αριθμητικά σε NaN
df['fixed acidity'] = pd.to_numeric(df['fixed acidity'], errors='coerce')
\end{lstlisting}

\subsubsection{Τι βρήκαμε}

\begin{table}[H]
\centering
\caption{Μη αριθμητικές τιμές ανά στήλη}
\label{tab:non_numeric}
\begin{tabular}{@{}lc@{}}
\toprule
\textbf{Μεταβλητή} & \textbf{Μη αριθμητικές τιμές} \\
\midrule
fixed acidity & 87 \\
volatile acidity & 45 \\
citric acid & 61 \\
residual sugar & 56 \\
chlorides & 57 \\
free sulfur dioxide & 56 \\
total sulfur dioxide & 53 \\
density & 53 \\
pH & 151 \\
sulphates & 63 \\
alcohol & 105 \\
quality & 58 \\
\midrule
\textbf{Σύνολο} & \textbf{845} \\
\bottomrule
\end{tabular}
\end{table}

\subsection{Κατηγορία 3: Αδύνατες/Ακραίες Τιμές}

\subsubsection{Τι ελέγξαμε}

Ακόμα και αν μια τιμή είναι αριθμητική, μπορεί να είναι \textbf{αδύνατη} ή \textbf{μη ρεαλιστική}. Ορίσαμε λογικά όρια για κάθε μεταβλητή βάσει επιστημονικής γνώσης:

\begin{table}[H]
\centering
\caption{Λογικά όρια τιμών ανά μεταβλητή}
\label{tab:limits}
\begin{tabular}{@{}lccp{5cm}@{}}
\toprule
\textbf{Μεταβλητή} & \textbf{Min} & \textbf{Max} & \textbf{Αιτιολόγηση} \\
\midrule
Fixed acidity & 0 & 20 & Τυπικές τιμές 4-16 g/L \\
Volatile acidity & 0 & 5 & Τυπικές τιμές 0.1-1.5 g/L \\
Citric acid & 0 & 3 & Τυπικές τιμές 0-1 g/L \\
Residual sugar & 0 & 100 & Ξηρά κρασιά <4, γλυκά ως 100 g/L \\
Chlorides & 0 & 1 & Τυπικές τιμές 0.01-0.2 g/L \\
Free SO\textsubscript{2} & 0 & 300 & Νομικό όριο ~300 mg/L \\
Total SO\textsubscript{2} & 0 & 500 & Νομικό όριο ~400 mg/L \\
Density & 0.9 & 1.1 & Κοντά στο νερό (~1 g/cm³) \\
pH & 2 & 5 & Κρασί: τυπικά 2.9-4.0 \\
Sulphates & 0 & 3 & Τυπικές τιμές 0.3-1.5 g/L \\
Alcohol & 5 & 20 & Κρασί: τυπικά 8-15\% \\
Quality & 0 & 10 & Κλίμακα βαθμολόγησης \\
\bottomrule
\end{tabular}
\end{table}

\subsubsection{Τι βρήκαμε}

\begin{table}[H]
\centering
\caption{Αδύνατες τιμές ανά μεταβλητή}
\label{tab:impossible}
\begin{tabular}{@{}lcp{6cm}@{}}
\toprule
\textbf{Μεταβλητή} & \textbf{Πλήθος} & \textbf{Παραδείγματα} \\
\midrule
fixed acidity & 156 & 999, -999, 9999 \\
volatile acidity & 170 & 999, -999, 9999 \\
citric acid & 164 & 999, -999, 9999 \\
residual sugar & 165 & 999, -999, 9999 \\
chlorides & 161 & 999, -999, 9999 \\
free sulfur dioxide & 151 & 999, -999, 9999 \\
total sulfur dioxide & 160 & 999, -999, 9999 \\
density & 154 & 999, -999, 9999 \\
pH & 202 & 0, 999, -999, 9999 \\
sulphates & 160 & 999, -999, 9999 \\
alcohol & 197 & 999, -999, 9999 \\
quality & 197 & 999, -999, 9999 \\
\midrule
\textbf{Σύνολο} & \textbf{2.037} & \\
\bottomrule
\end{tabular}
\end{table}

\begin{warningbox}[Παρατήρηση]
Οι τιμές 999, -999 και 9999 εμφανίζονται συστηματικά σε όλες τις στήλες. Αυτό υποδηλώνει ότι χρησιμοποιήθηκαν \textbf{ως κωδικοί για ελλείπουσες τιμές} κατά την καταχώρηση των δεδομένων, αντί να αφεθούν κενά τα πεδία.
\end{warningbox}

\subsection{Κατηγορία 4: Διπλοεγγραφές (Duplicates)}

\subsubsection{Τι είναι οι διπλοεγγραφές}

Μια διπλοεγγραφή είναι μια γραμμή που εμφανίζεται \textbf{πανομοιότυπη} με άλλη γραμμή στο dataset. Αυτό μπορεί να συμβεί από:
\begin{itemize}
    \item Λάθος κατά την εισαγωγή δεδομένων
    \item Συγχώνευση πολλαπλών αρχείων χωρίς έλεγχο
    \item Τεχνικά προβλήματα κατά την αποθήκευση
\end{itemize}

\subsubsection{Τι βρήκαμε}

\begin{lstlisting}[caption={Εντοπισμός διπλοεγγραφών}]
# Εύρεση διπλοεγγραφών
duplicates = df.duplicated()
print(f"Πλήθος διπλοεγγραφών: {duplicates.sum()}")
\end{lstlisting}

\begin{successbox}[Αποτέλεσμα]
Εντοπίστηκαν \textbf{1.434 διπλοεγγραφές} στο δείγμα.
\end{successbox}

\subsection{Κατηγορία 5: Ελλείπουσες Τιμές (Missing Values)}

\subsubsection{Τι είναι οι ελλείπουσες τιμές}

\begin{infobox}[Ορισμός: Missing Values]
Οι ελλείπουσες τιμές (ή NaN - Not a Number) είναι θέσεις στο dataset όπου δεν υπάρχει καταχωρημένη τιμή. Μπορεί να προκύψουν από:
\begin{itemize}
    \item Αδυναμία μέτρησης ή καταγραφής
    \item Λάθη κατά τη μεταφορά δεδομένων
    \item Μετατροπή μη έγκυρων τιμών σε NaN (όπως στη δική μας περίπτωση)
\end{itemize}
\end{infobox}

\subsubsection{Τι βρήκαμε}

Μετά τη μετατροπή των μη έγκυρων τιμών σε NaN, έχουμε:

\begin{table}[H]
\centering
\caption{Ελλείπουσες τιμές ανά μεταβλητή}
\label{tab:missing}
\begin{tabular}{@{}lcc@{}}
\toprule
\textbf{Μεταβλητή} & \textbf{Missing Values} & \textbf{Ποσοστό} \\
\midrule
fixed acidity & 574 & 10.88\% \\
volatile acidity & 542 & 10.28\% \\
citric acid & 1 & 0.02\% \\
residual sugar & 134 & 2.54\% \\
chlorides & 8 & 0.15\% \\
free sulfur dioxide & 1 & 0.02\% \\
total sulfur dioxide & 1 & 0.02\% \\
density & 545 & 10.33\% \\
pH & 644 & 12.21\% \\
sulphates & 1 & 0.02\% \\
alcohol & 600 & 11.38\% \\
quality & 40 & 0.76\% \\
wine\_type & 1 & 0.02\% \\
\midrule
\textbf{Σύνολο} & \textbf{3.092} & \\
\bottomrule
\end{tabular}
\end{table}

% ==============================================================================
\section{Αντιμετώπιση Προβλημάτων}
% ==============================================================================

\subsection{Στρατηγική Αντιμετώπισης}

Για κάθε κατηγορία προβλήματος, εφαρμόσαμε κατάλληλη στρατηγική:

\begin{table}[H]
\centering
\caption{Στρατηγική αντιμετώπισης ανά τύπο προβλήματος}
\label{tab:strategy}
\begin{tabular}{@{}lp{10cm}@{}}
\toprule
\textbf{Πρόβλημα} & \textbf{Στρατηγική Αντιμετώπισης} \\
\midrule
Τυπογραφικά λάθη & Διόρθωση στη σωστή τιμή ή μετατροπή σε NaN \\
Μη αριθμητικές τιμές & Μετατροπή σε NaN \\
Αδύνατες τιμές & Μετατροπή σε NaN \\
Διπλοεγγραφές & Αφαίρεση (κράτηση μόνο της πρώτης εμφάνισης) \\
Missing values & Συμπλήρωση με τη διάμεσο (median) ανά τύπο κρασιού \\
\bottomrule
\end{tabular}
\end{table}

\subsection{Διόρθωση Τυπογραφικών Λαθών}

\begin{lstlisting}[caption={Διόρθωση τυπογραφικών λαθών}]
# Ορισμός αντιστοιχίσεων διόρθωσης
wine_type_corrections = {
    'Redd ': 'red',     # Τυπογραφικό λάθος
    'r3d': 'red',       # Τυπογραφικό λάθος (αριθμός 3 αντί e)
    'WHITTE': 'white',  # Τυπογραφικό λάθος
    'whit': 'white',    # Ελλιπής λέξη
    'whate': 'white',   # Τυπογραφικό λάθος
    '999': np.nan,      # Κωδικός ελλείπουσας τιμής
    '9999': np.nan,     # Κωδικός ελλείπουσας τιμής
    '-999': np.nan,     # Κωδικός ελλείπουσας τιμής
    '?': np.nan,        # Σύμβολο αγνώστου
}

# Εφαρμογή διορθώσεων
for wrong, correct in wine_type_corrections.items():
    mask = df['wine_type'] == wrong
    df.loc[mask, 'wine_type'] = correct
\end{lstlisting}

\begin{successbox}[Αποτέλεσμα]
Διορθώθηκαν \textbf{560 εγγραφές} με τυπογραφικά λάθη στη στήλη wine\_type.
\end{successbox}

\subsection{Αντιμετώπιση Αδύνατων Τιμών}

\begin{lstlisting}[caption={Αντικατάσταση αδύνατων τιμών}]
# Ορισμός ορίων για κάθε μεταβλητή
value_limits = {
    'fixed acidity': (0, 20),
    'volatile acidity': (0, 5),
    'pH': (2, 5),
    # ... κλπ
}

# Για κάθε μεταβλητή
for col, (min_val, max_val) in value_limits.items():
    # Εντοπισμός τιμών εκτός ορίων
    out_of_range = (df[col] < min_val) | (df[col] > max_val)
    # Αντικατάσταση με NaN
    df.loc[out_of_range, col] = np.nan
\end{lstlisting}

\subsection{Αφαίρεση Διπλοεγγραφών}

\begin{lstlisting}[caption={Αφαίρεση διπλοεγγραφών}]
# Αφαίρεση διπλοεγγραφών, κράτηση πρώτης εμφάνισης
df = df.drop_duplicates()
df = df.reset_index(drop=True)
\end{lstlisting}

\begin{successbox}[Αποτέλεσμα]
Αφαιρέθηκαν \textbf{1.434 διπλοεγγραφές}.
\end{successbox}

\subsection{Συμπλήρωση Ελλειπουσών Τιμών}

\subsubsection{Επιλογή Μεθόδου: Διάμεσος (Median)}

\begin{infobox}[Γιατί Διάμεσος και όχι Μέσος Όρος;]
\begin{itemize}
    \item Η \textbf{διάμεσος} είναι η μεσαία τιμή όταν τα δεδομένα ταξινομηθούν
    \item Είναι \textbf{ανθεκτική} σε ακραίες τιμές (robust)
    \item Ο \textbf{μέσος όρος} επηρεάζεται πολύ από ακραίες τιμές
    \item Για δεδομένα με ασυμμετρία (skewness), η διάμεσος είναι καταλληλότερη
\end{itemize}
\end{infobox}

\subsubsection{Γιατί ανά τύπο κρασιού;}

Τα κόκκινα και λευκά κρασιά έχουν \textbf{διαφορετικά χημικά χαρακτηριστικά}:

\begin{table}[H]
\centering
\caption{Σύγκριση διαμέσων τιμών ανά τύπο κρασιού}
\begin{tabular}{@{}lcc@{}}
\toprule
\textbf{Μεταβλητή} & \textbf{Κόκκινο} & \textbf{Λευκό} \\
\midrule
Fixed acidity & 7.80 & 6.80 \\
Volatile acidity & 0.50 & 0.26 \\
Total sulfur dioxide & ~40 & ~133 \\
Alcohol & 10.20 & 10.40 \\
\bottomrule
\end{tabular}
\end{table}

Αν χρησιμοποιούσαμε τη συνολική διάμεσο, θα εισάγαμε σφάλμα στις εκτιμήσεις.

\subsubsection{Υλοποίηση}

\begin{lstlisting}[caption={Συμπλήρωση missing values με διάμεσο}]
# Για κάθε αριθμητική μεταβλητή
for col in numeric_columns:
    # Για κάθε τύπο κρασιού
    for wine_type in ['red', 'white']:
        # Βρες τις γραμμές με missing value
        mask = (df['wine_type'] == wine_type) & (df[col].isnull())
        
        # Υπολόγισε τη διάμεσο για τον συγκεκριμένο τύπο
        median_val = df.loc[df['wine_type'] == wine_type, col].median()
        
        # Συμπλήρωσε τα missing με τη διάμεσο
        df.loc[mask, col] = median_val

# Για τη μεταβλητή wine_type, χρήση mode (πιο συχνή τιμή)
mode_wine = df['wine_type'].mode()[0]  # 'white'
df['wine_type'] = df['wine_type'].fillna(mode_wine)
\end{lstlisting}

% ==============================================================================
\section{Συνοπτικά Αποτελέσματα}
% ==============================================================================

\subsection{Πίνακας Προβλημάτων και Αντιμετώπισης}

\begin{table}[H]
\centering
\caption{Συνοπτικός πίνακας προβλημάτων}
\label{tab:summary}
\begin{tabular}{@{}lcc@{}}
\toprule
\textbf{Τύπος Προβλήματος} & \textbf{Πλήθος} & \textbf{Αντιμετώπιση} \\
\midrule
Περιττές στήλες (Unnamed) & 2 & Αφαίρεση \\
Τυπογραφικά λάθη wine\_type & 9 τύποι (560) & Διόρθωση/NaN \\
Μη αριθμητικές τιμές & 845 & Μετατροπή σε NaN \\
Αδύνατες/ακραίες τιμές & 2.037 & Μετατροπή σε NaN \\
Διπλοεγγραφές & 1.434 & Αφαίρεση \\
Missing values (αρχικά) & 3.092 & Συμπλήρωση με διάμεσο \\
Missing values (τελικά) & 0 & -- \\
\bottomrule
\end{tabular}
\end{table}

\subsection{Μεταβολή Μεγέθους Dataset}

\begin{figure}[H]
\centering
\begin{tabular}{|c|c|c|}
\hline
\textbf{Στάδιο} & \textbf{Γραμμές} & \textbf{Στήλες} \\
\hline
Αρχικό dataset & 8.385 & 15 \\
\hline
Μετά δειγματοληψία (80\%) & 6.708 & 13 \\
\hline
Μετά αφαίρεση διπλοεγγραφών & 5.274 & 13 \\
\hline
\textbf{Τελικό καθαρό dataset} & \textbf{5.274} & \textbf{13} \\
\hline
\end{tabular}
\caption{Εξέλιξη μεγέθους dataset}
\end{figure}

\subsection{Κατανομή Τελικού Dataset}

\begin{table}[H]
\centering
\caption{Κατανομή τελικού dataset ανά τύπο κρασιού}
\begin{tabular}{@{}lcc@{}}
\toprule
\textbf{Τύπος Κρασιού} & \textbf{Πλήθος} & \textbf{Ποσοστό} \\
\midrule
Λευκό (white) & 3.924 & 74.4\% \\
Κόκκινο (red) & 1.350 & 25.6\% \\
\midrule
\textbf{Σύνολο} & \textbf{5.274} & \textbf{100\%} \\
\bottomrule
\end{tabular}
\end{table}

% ==============================================================================
\section{Συμπεράσματα}
% ==============================================================================

\subsection{Τι Μάθαμε}

\begin{enumerate}
    \item Τα πραγματικά δεδομένα περιέχουν \textbf{πολλαπλά είδη προβλημάτων}
    \item Ο καθαρισμός δεδομένων είναι \textbf{απαραίτητο προκαταρκτικό βήμα}
    \item Κάθε τύπος προβλήματος απαιτεί \textbf{διαφορετική αντιμετώπιση}
    \item Οι αποφάσεις καθαρισμού πρέπει να είναι \textbf{τεκμηριωμένες και αναπαραγώγιμες}
\end{enumerate}

\subsection{Περιορισμοί}

\begin{itemize}
    \item Η συμπλήρωση με διάμεσο \textbf{μειώνει τη μεταβλητότητα} των δεδομένων
    \item Η αφαίρεση διπλοεγγραφών \textbf{μείωσε το μέγεθος} του dataset κατά 21\%
    \item Ορισμένες αποφάσεις (π.χ. όρια τιμών) βασίστηκαν σε \textbf{εμπειρική γνώση}
\end{itemize}

\subsection{Επόμενα Βήματα}

Το καθαρισμένο dataset είναι πλέον έτοιμο για:
\begin{itemize}
    \item Περιγραφική στατιστική ανάλυση
    \item Οπτικοποίηση (boxplots, histograms, κλπ.)
    \item Εντοπισμό ακραίων τιμών (outliers)
    \item Κατασκευή προβλεπτικών μοντέλων
\end{itemize}

% ==============================================================================
\appendix
\section{Παράρτημα: Πλήρης Κώδικας Python}
% ==============================================================================

\begin{lstlisting}[caption={Πλήρης κώδικας καθαρισμού δεδομένων}]
import pandas as pd
import numpy as np

# Φόρτωση δεδομένων
df_original = pd.read_csv('Data Analysis_2026 1st Case_Data.csv')

# Αφαίρεση περιττών στηλών
unnamed_cols = [col for col in df_original.columns if 'Unnamed' in col]
df_original = df_original.drop(columns=unnamed_cols)

# Δειγματοληψία
AEM = 6931
df = df_original.sample(frac=0.80, random_state=AEM)
df = df.reset_index(drop=True)

# Διόρθωση τυπογραφικών λαθών
wine_type_corrections = {
    'Redd ': 'red', 'r3d': 'red',
    'WHITTE': 'white', 'whit': 'white', 'whate': 'white',
    '999': np.nan, '9999': np.nan, '-999': np.nan, '?': np.nan
}
for wrong, correct in wine_type_corrections.items():
    df.loc[df['wine_type'] == wrong, 'wine_type'] = correct

# Μετατροπή σε αριθμητικές τιμές
numeric_columns = ['fixed acidity', 'volatile acidity', 'citric acid',
                   'residual sugar', 'chlorides', 'free sulfur dioxide',
                   'total sulfur dioxide', 'density', 'pH', 'sulphates',
                   'alcohol', 'quality']
for col in numeric_columns:
    df[col] = pd.to_numeric(df[col], errors='coerce')

# Αντικατάσταση αδύνατων τιμών
value_limits = {
    'fixed acidity': (0, 20), 'volatile acidity': (0, 5),
    'citric acid': (0, 3), 'residual sugar': (0, 100),
    'chlorides': (0, 1), 'free sulfur dioxide': (0, 300),
    'total sulfur dioxide': (0, 500), 'density': (0.9, 1.1),
    'pH': (2, 5), 'sulphates': (0, 3), 'alcohol': (5, 20),
    'quality': (0, 10)
}
for col, (min_val, max_val) in value_limits.items():
    mask = (df[col] < min_val) | (df[col] > max_val)
    df.loc[mask, col] = np.nan

# Αφαίρεση διπλοεγγραφών
df = df.drop_duplicates().reset_index(drop=True)

# Συμπλήρωση missing values
df['wine_type'] = df['wine_type'].fillna(df['wine_type'].mode()[0])
for col in numeric_columns:
    for wine_type in ['red', 'white']:
        mask = (df['wine_type'] == wine_type) & (df[col].isnull())
        median = df.loc[df['wine_type'] == wine_type, col].median()
        df.loc[mask, col] = median

# Αποθήκευση
df.to_csv('cleaned_wine_data_AEM_6931.csv', index=False)
\end{lstlisting}

\end{document}
